% !TeX root = ../main.tex
\section{Read/Write Mix}\label{sec:rw-mix}

\begin{hypothesis}
  Every written line needs a read-for-ownership (RFO) and later a write-back, so it costs about two lines
  of DRAM traffic plus bus turnarounds. Demand bandwidth should therefore fall as the write fraction grows
  while DRAM traffic stays nearly constant, and loaded latency should be highest for write-heavy mixes,
  whose write-backs compete with the probe's reads.
\end{hypothesis}

\begin{knobs}
  Fixed       & \qty{128}{\mebi\byte} per generator core; sequential; ordinary (write-allocate) full-line stores; \qty{4}{\kibi\byte} pages \\
  Varied      & reads per 10 \emph{lines}: 10 (100/0), 7 (70/30), 5 (50/50), 0 (0/100); $T = 1$ and $9$ generator cores \\
  Latency     & concurrent random pointer-chase probe on CPU~2 (\code{mem_loaded}); $T=0$ gives the unloaded value \\
  Script      & \code{Cache_and_Memory_Performance_Profiling/scripts/run.sh rw_mix} \\
\end{knobs}

\paragraph{Mix definition.} Of every 10 consecutive \qty{64}{\byte} lines, the first $R$ are read and
the remaining $10-R$ are fully overwritten, so the ratio holds for the traffic the memory system sees.
Mixing \qty{8}{\byte} elements instead is a trap: every source line is read \emph{and} every destination
line written, so a nominal ``70/30'' moves about three lines per useful line (\cref{sec:ai-use}).

\begin{figure}[tbp]
  \centering
  \begin{subfigure}[t]{0.49\linewidth}
    \IfDataExists{rw_mix}{%
    \begin{tikzpicture}
      \begin{axis}[acs, height=0.72\linewidth, ymin=0, ymax=50, xtick={0,5,7,10},
                   xticklabels={0/100,50/50,70/30,100/0}, xlabel={Read/write mix (lines)},
                   ylabel={Bandwidth (GiB/s)},
                   legend style={font=\scriptsize, inner sep=1pt, at={(0.5,1.03)}, anchor=south},
                   legend columns=2]
        \addplot+[blue, mark=*, mark options={fill=blue}, iqr bars] table[x=reads10, y=total_gib_s_median, iqr of=total_gib_s, only if={threads}{9}] {data/rw_mix.csv};
        \addlegendentry{demand, $T=9$}
        \addplot+[blue, mark=o, dashed] table[x=reads10, y=traffic_gib_s_median, only if={threads}{9}] {data/rw_mix.csv};
        \addlegendentry{est.\ traffic, $T=9$}
        \addplot+[blue!50!black, only marks, mark=x, mark size=3pt] table[x=reads10, y expr=\thisrow{imc_rd_gib_s_median}+\thisrow{imc_wr_gib_s_median}, only if={threads}{9}] {data/rw_mix.csv};
        \addlegendentry{DRAM (IMC), $T=9$}
        \addplot+[red, mark=square*, mark options={fill=red}, iqr bars] table[x=reads10, y=total_gib_s_median, iqr of=total_gib_s, only if={threads}{1}] {data/rw_mix.csv};
        \addlegendentry{demand, $T=1$}
        \addplot+[red, mark=square, dashed] table[x=reads10, y=traffic_gib_s_median, only if={threads}{1}] {data/rw_mix.csv};
        \addlegendentry{est.\ traffic, $T=1$}
        \addplot+[red!50!black, only marks, mark=+, mark size=3pt] table[x=reads10, y expr=\thisrow{imc_rd_gib_s_median}+\thisrow{imc_wr_gib_s_median}, only if={threads}{1}] {data/rw_mix.csv};
        \addlegendentry{DRAM (IMC), $T=1$}
      \end{axis}
    \end{tikzpicture}}
    \caption{Throughput.}
  \end{subfigure}\hfill
  \begin{subfigure}[t]{0.49\linewidth}
    \IfDataExists{rw_mix}{%
    \begin{tikzpicture}
      \begin{axis}[acs, height=0.72\linewidth, ymin=0, xtick={0,5,7,10},
                   xticklabels={0/100,50/50,70/30,100/0}, xlabel={Read/write mix (lines)},
                   ylabel={Probe latency (ns)}, legend pos=north east, legend style={font=\scriptsize}]
        \addplot+[iqr bars] table[x=reads10, y=probe_ns_median, iqr of=probe_ns, only if={threads}{9}] {data/rw_mix.csv};
        \addlegendentry{$T=9$}
        \addplot+[iqr bars] table[x=reads10, y=probe_ns_median, iqr of=probe_ns, only if={threads}{1}] {data/rw_mix.csv};
        \addlegendentry{$T=1$}
      \end{axis}
    \end{tikzpicture}}
    \caption{Loaded latency of a random read.}
  \end{subfigure}
  \caption{Read/write mixes: throughput and concurrently measured latency. Demand: bytes the program reads
  and writes; est.\ traffic: adds one RFO line per written line; DRAM (IMC): memory-controller read+write
  count over the same windows. The unloaded probe latency
  is \valqty{rw_unloaded_ns}{\nano\second}.}\label{fig:rw-mix}
\end{figure}

\subsection{Analysis}
With nine generator cores, demand bandwidth is \valqty{rw_10_t9_gib}{\gibi\byte\per\second} for 100\%
reads, \valqty{rw_7_t9_gib}{\gibi\byte\per\second} at 70/30, \valqty{rw_5_t9_gib}{\gibi\byte\per\second}
at 50/50 and \valqty{rw_0_t9_gib}{\gibi\byte\per\second} for 100\% writes, falling with the write fraction as
predicted. The memory-controller counters measure \valqty{rw_10_t9_imc}{\gibi\byte\per\second} to
\valqty{rw_0_t9_imc}{\gibi\byte\per\second} across the mixes, within about 1\% of the RFO-counted estimate
(\valqty{rw_10_t9_traffic}{\gibi\byte\per\second} to \valqty{rw_0_t9_traffic}{\gibi\byte\per\second}; idle
background \valqty{cc_imc_rd_idle}{\gibi\byte\per\second}); for pure writes the controller reads
\valqty{rw_0_t9_imcrd}{\gibi\byte\per\second} (the RFOs) and writes \valqty{rw_0_t9_imcwr}{\gibi\byte\per\second}.
Total DRAM traffic stays at about \qty{30}{\gibi\byte\per\second} whatever the mix, as does MLC's
peak-injection bandwidth (\valqty{mlc_pk_allreads_gib}{\gibi\byte\per\second} all reads,
\valqty{mlc_pk_r31_gib}{\gibi\byte\per\second} at 3:1, \valqty{mlc_pk_r11_gib}{\gibi\byte\per\second} at 1:1).
The cost of writes is the RFO, not a loss of controller efficiency.

Probe latency rises from \valqty{rw_10_t9_ns}{\nano\second} (reads) to \valqty{rw_0_t9_ns}{\nano\second}
(writes): at the same total DRAM traffic, pure writes double it (\val{rw_0_t9_ns_rel}$\times$) and the 70/30 and 50/50 mixes raise it by about a third
(\val{rw_7_t9_ns_rel}$\times$, \val{rw_5_t9_ns_rel}$\times$): write-backs and RFOs share the probe's queues,
and bus turnarounds lengthen every read's wait.

One generator core ($T=1$) is not controller-bound: demand bandwidth falls only from
\valqty{rw_10_t1_gib}{\gibi\byte\per\second} to \valqty{rw_0_t1_gib}{\gibi\byte\per\second} (perhaps because
stores retire without waiting for their RFO; untested), while DRAM traffic grows from
\valqty{rw_10_t1_imc}{\gibi\byte\per\second} to \valqty{rw_0_t1_imc}{\gibi\byte\per\second}. Here the counters
read \val{rw_t1_imcest_min}--\val{rw_t1_imcest_max}$\times$ the RFO-counted estimate, even for pure reads, a
gap we did not resolve; \emph{untested hypothesis:} the estimate counts lines still in the L3 from the
previous pass, which a buffer many times the L3 would test. The probe barely notices one core
(\val{rw_t1_ns_min}--\valqty{rw_t1_ns_max}{\nano\second}). On all ten cores, \code{copy} delivers
\valqty{rw_ref_copy}{\gibi\byte\per\second} and \code{triad} \valqty{rw_ref_triad}{\gibi\byte\per\second}; copy reads
\valqty{rw_ref_copy_imcrd}{\gibi\byte\per\second} and writes \valqty{rw_ref_copy_imcwr}{\gibi\byte\per\second}
(2:1, source plus RFO), and triad reads \valqty{rw_ref_triad_imcrd}{\gibi\byte\per\second} and writes
\valqty{rw_ref_triad_imcwr}{\gibi\byte\per\second} (3:1). Non-temporal stores (no RFO) were out of scope.
